
% Retained as checking evidence: maps:src:E.1.3
\begin{theorem}[the integral cusp calculation]\hypertarget{tag-02M6}{}\label{maps:src:E.1.3}
\StatementRecord{maps:src:E.1.3}{proved}{02M6}{\eqref{maps:eq:E.5}, \eqref{maps:eq:E.7}; \cite[Proposition~8.3.5, pp.~42--43]{mapsDeligne1974HodgeIII}; \cite[Theorem B]{ffChen}}


For the central fibre $W_3$ of \eqref{maps:eq:E.5},
\[
 H_q(W_3;\mathbb Z)=\mathbb Z^{b_q},\qquad
 (b_0,\ldots,b_6)=(1,3,9,9,10,3,1).
 \tag{E.6}\hypertarget{tag-029T}{}\label{maps:eq:E.6}
\]
Integral specialization is onto the monodromy invariant lattice in
every degree.  Its only kernel is a copy of $\mathbb Z$ in degree four.

\end{theorem}


\begin{proof}
\begin{enumerate}[label={$\langle1\rangle$\arabic*.},leftmargin=*]
\item The normalization is the three-dimensional permutohedral toric variety.
Its paired boundary surfaces consist of four copies of $\mathrm{dP}_6$
and three copies of $\mathbb P^1\times\mathbb P^1$.  There are twelve
normalized triple curves and six quadruple points.  The normalization
complex for the constant sheaf is the alternating branch complex.
Local exactness follows from the augmented simplex on the branches
of $z_0z_1z_2z_3=0$.  Its nonzero rows after taking cohomology are
\[
\begin{array}{c|rrrr|rr}
q& p=0&p=1&p=2&p=3&\operatorname{rank}d^0&
 \operatorname{rank}d^1\ \text{or }d^2\\ \hline
0&1&7&12&6&0&4,5\\
2&11&22&12&0&5&9\\
4&11&7&0&0&4&0\\
6&1&0&0&0&0&0
\end{array}
\tag{E.7}\hypertarget{tag-029U}{}\label{maps:eq:E.7}
\]
Every nonzero Smith invariant in these matrices is one.

\item The integral matrices are as follows.
The fourteen rays are $\pm v$, with
$v\in\{0,1\}^3\setminus\{0\}$.  The twenty-four maximal cones
are the translated Kuhn tetrahedra containing zero.  This enumeration
requires exactly the eight cubes with lower vertex in $\{-1,0\}^3$;
no omitted cube can contain zero.  The Picard group of the normalization
is the group of ray coefficients modulo the three character relations.
We set the coefficients at $e_1,e_2,e_3$ to zero, giving an integral
basis of rank eleven.  On a boundary surface $D_v$, choose adjacent
rays $u_1,u_2$ for which $(v,u_1,u_2)$ is unimodular and set the two
remaining character coefficients to zero.  Restriction is evaluated
by subtracting the unique integral linear function with the prescribed
values at these three rays.  The paired restriction is
$\operatorname{res}_{D_v}-\operatorname{res}_{D_{-v}}$, with the
second adjacent ray corresponding to $u-v$.

\item For a triangle $(0,v,u)$ with $0<v<u$ coordinatewise, the next map
is the alternating sum of the degrees on the three boundary curves
of $D_{u-v},D_u,D_v$.  In a surface fan, if the neighbours of a ray
$w$ are $h_1,h_2$ and $h_1+h_2=bw$ modulo $v$, that degree is
$a_{h_1}+a_{h_2}-ba_w$.  In degree four, the restriction map is dual
to the seven divisor differences $[D_v]-[D_{-v}]$.  In degree zero
the matrices are the cellular incidence matrices of the one-vertex
Kuhn triangulation of $T^3$.  These rules define \eqref{maps:eq:E.7} over
$\mathbb Z$, and both compositions are zero.
Its exact Smith reductions give the ranks and primitivity just stated.

\item For completeness, the possible higher differential of this
normalization spectral sequence vanishes integrally.  Replace $C_3$
by $hC_3$, $0\leq h\leq1$.  The preceding uniform properness estimate
applies simultaneously.  Local parameter vector fields preserving
$t_c$ and each coordinate stratum patch by a partition of unity;
their flow identifies the central fibres and their specialization
maps.  We may therefore use $C_3=0$.  This toroidal degeneration is
projective after passing to the finite cover with deck lattice
$2\mathbb Z^3$.  An explicit coherent polarization is given by
\[
 Q(n)=4\sum n_i^2-(\sum n_i)^2.
\]
For the tetrahedron ordered $1,2,3$, its affine interpolant is
$3n_1+n_2-n_3$, and their difference is
\[
 \sum_i n_i(n_i-1)+
 \sum_{i<j}(n_i-n_j)(n_i-n_j-1)\geq0
 \quad(n\in\mathbb Z^3).
\]
Equality holds exactly at its four vertices.  Translation and
permutation give strict coherence on every face.  The finite cover
has simple normal crossings, and its descent complex is a complex
of pure toric cohomology groups.  Over $\mathbb Q$ the higher
differentials vanish by their different weights \cite[Proposition~8.3.5, pp.~42--43]{mapsDeligne1974HodgeIII}.  Taking finite
deck invariants returns the original fibre.  The integral $E_2$
groups of \eqref{maps:eq:E.7} are free, so a higher differential which is zero
over $\mathbb Q$ is zero integrally.  Summing its diagonals proves
\eqref{maps:eq:E.6}, with no extension torsion.

\item The same projective comparison allows the local invariant cycle
theorem, in the form of \cite[Theorem B]{ffChen}, followed by exact
finite-group invariants over $\mathbb Q$.  The cusp monodromy on
$H_1$ is $\left(\begin{smallmatrix}I_3&I_3\\0&I_3\end{smallmatrix}\right)$.
Its exterior-power invariant dimensions are
$(1,3,9,9,9,3,1)$, by a computation with exact exterior matrices.  Comparison with \eqref{maps:eq:E.6} proves the last assertion.
We now strengthen this to an integral statement.  Write $b_i$ for the
three degree-one invariant deck classes and $a_i$ for the angular
classes, so the infinitesimal monodromy sends $a_i$ to $b_i$.
At $C_3=0$, every integral symmetric matrix $S$ gives a line bundle
on the toroidal family.  Its descent data are as follows.  Assign to a fan vertex $n$
the integer
\[
 q_S(n)=\frac{n^tSn-\sum_iS_{ii}n_i}{2}.
\]
Interpolate on each unimodular Kuhn tetrahedron.  This is an integral
piecewise linear support function.  Translation by $\lambda$ changes
it by the integral affine function
$n^tS\lambda+q_S(\lambda)$, giving the deck linearization
$t_c^{-q_S(\lambda)}x^{-S\lambda}$.
The identity
$q_S(\lambda+\mu)=q_S(\lambda)+q_S(\mu)+\lambda^tS\mu$
is exactly the cocycle condition.  On a smooth fibre its Chern class
is $\sum S_{ij}a_i b_j$.  Restriction of this line bundle to the
central fibre therefore supplies every class
\[
 a_i b_i,\qquad a_i b_j+a_j b_i\quad(i<j)
\]
in the specialization image, as well as the $b_i$ and their products.
These elements generate the full integral invariant exterior algebra
for three pairs.  To check saturation, one enumerates all monomials of weighted degree at most six in these
nine generators and reduces their coordinates in an integral kernel
basis of $\bigwedge^q U-I$.  In every degree the nonzero Smith factors
are all one.  This finite check has only $2^3$ possible products of
the odd generators and at most three even factors.  It can also be
seen directly: in degree three the mixed kernel is the trace-zero
rank-eight lattice, generated by six off-diagonal terms and two
primitive diagonal differences; the complementary degrees are
generated by the corresponding products of two or three $a_i b_i$.
The parameter isotopy already constructed identifies the integral
specialization maps for general $C_3$.  Thus the images are primitive
and surjective in every degree.  Comparing their ranks with \eqref{maps:eq:E.6}
gives exactly the asserted free degree-four kernel.
\end{enumerate}
\end{proof}

% Retained as checking evidence: maps:src:F.2.5 rational and characteristic re-proof
\begin{theorem}[The integral eight-manifold]\hypertarget{tag-02NW}{}\label{maps:src:F.2.5}
\StatementRecord{maps:src:F.2.5}{proved}{02NW}{\SeriesRef{maps:src:F.2.1}, \SeriesRef{night:P1.1}, \SeriesRef{night:P1.9}, \SeriesRef{night:P1.11}; \cite{ffAchingerOgus20,StacksProject,HirzebruchTopologicalMethods,ffPetersSterk}}


There is a closed smooth oriented eight-manifold with the prescribed
index-one Kuhn cusp, two free affine fillings of multiplicities $5,4$,
trivial fundamental group, and rational Betti numbers
\[
 (1,0,0,0,4,0,0,0,1). \tag{F.5}\hypertarget{tag-02F9}{}\label{maps:eq:F.5}
\]
It admits an almost complex structure.  Its signature is two and
its middle intersection lattice modulo torsion is
$3\langle1\rangle\oplus\langle-1\rangle$.
Its integral homology is computed in \cref{maps:src:F.2.6}, and a complex structure on it is constructed in \cref{maps:src:F.3.7}.

\end{theorem}


\begin{proof}
\begin{enumerate}[label={$\langle1\rangle$\arabic*.}]
\item Use the matrices and affine vectors of \cref{pub:V:rank-six-data}.
Theorem~\ref{night:P1.1} gives their orders, the primitive cusp
marking, the coinvariant homomorphism $\ell(x)=x_0$, and freeness of
every nonidentity power of the affine actions. Consequently
\[
 (T^6\times\Delta_u)/
 \langle(x,u)\mapsto(M_ix+v_i/m_i,\zeta_{m_i}u)\rangle,
 \qquad (m_1,m_2)=(5,4),
\]
are smooth oriented finite caps. On the universal angular cover
$u=r e^{2\pi ia}$ their marked collar maps are
\[
 \alpha_i(x,u)=x-av_i,
 \qquad
 \alpha_i(M_ix+v_i/m_i,\zeta_{m_i}u)=M_i\alpha_i(x,u).
\]
At the cusp we take the index-one Kuhn cap of
Theorem~\ref{night:P1.9}, in the marking of
Proposition~\ref{maps:src:F.2.1}. The three displayed boundary bundles
therefore glue to the marked torus bundle over the intervening pair
of pants, with zero additional clutch translation. Every boundary
is paired. The result is a closed connected smooth eight-manifold,
oriented by the cusp complex structure and the continued fibre/base
orientation.

\item The small cusp cap retracts onto its central fibre $C$ through
the mapping-cylinder homeomorphism
\cite[Theorem 4.1.1 and Lemmas 4.1.2--4.1.3, pp.~480--482]{ffAchingerOgus20}.
Indeed its divisorial log chart is
$\mathbb N\to\mathbb N^r$, $1\mapsto(1,\ldots,1)$; it is log smooth,
saturated and vertical, and the cap is proper by
Theorem~\ref{night:P1.9}. In the resulting cylinder the homotopy
$[x,r]\mapsto[x,(1-t)r]$ fixes $C$ and retracts onto it.

Before the period quotient, the components of $C$ and their
nonempty intersections are smooth complete toric strata. Each is
simply connected: a loop moves into its dense torus by general
position, and the one-parameter loops along the rays of a smooth
maximal cone bound the coordinate discs; those rays are an integral
basis. Van Kampen sends component-crossing paths to edges and triple
intersections to triangles of the Kuhn nerve in $\mathbb R^3$.
Thus the covering core has trivial fundamental group. Its period
quotient has deck group $\mathbb Z^3$, so $\pi_1(C)=\mathbb Z^3$.
The nearby-torus map is the quotient of $\Lambda$ by the saturated
image of $M_0-I$, by the $H_1$ collapse in
Theorem~\ref{night:P1.11}. It is onto, and the cusp meridian dies
because its zero section extends.

We also record the invariant normal closure needed below. Let
$K$ be the smallest $\langle M_1,M_2\rangle$-invariant subgroup
containing $(M_0-I)\Lambda$, and let
$J=(M_1-I)\Lambda+(M_2-I)\Lambda=\ker\ell$.
The equality $\ell M_i=\ell$ gives $K\subseteq\ker\ell$.
On the abelian quotient $\Lambda/K$, the induced maps $a,b$
satisfy $ab=1$, $a^5=1$ and $b^4=1$. Thus $a^4=a^5=1$, so
$a=b=1$ and $J\subseteq K$. Consequently $K=J=\ker\ell$ as
integral subgroups.


Let $a_i$ be the two finite peripheral lifts. The universal covering
deck transformation of the $i$th cap has $m_i$th power equal to
translation by $v_i$, since $M_iv_i=v_i$. The collar map in
Step~1 sends $a_i$ to that transformation. Van Kampen therefore gives
\[
 a_i\lambda a_i^{-1}=M_i\lambda,
 \qquad a_1^5=v_1,\quad a_2^4=v_2,\quad a_1a_2=1.
\]
Write $a_1=t$ and $a_2=t^{-1}$. Then
$t=t^5t^{-4}=v_1+v_2$ lies in the abelian fibre subgroup. Conjugation
is consequently trivial, and that subgroup is a quotient of
$\Lambda_{\langle M_1,M_2\rangle}=\mathbb Z$, with coordinate $\ell$.
Now $\ell(v_1+v_2)=0$, so $t=0$, while $t^5=v_1$ kills its generator
$\ell(v_1)=1$. Every generator vanishes, proving $\pi_1=1$.
Equivalently, after this abelian reduction, the relations on
$(\ell,a_1,a_2)$ have matrix
\[
 \begin{pmatrix}-1&5&0\\1&0&4\\0&1&1\end{pmatrix},
 \qquad\det=-1.
\]

\item We calculate rational cohomology directly from the marked
restriction maps. For $1\leq q\leq4$ put
\[
 L_q=\bigwedge^q\Lambda^\vee\otimes\mathbb Q,
 \quad A_q=\bigwedge^qA,\quad B_q=\bigwedge^qB,
 \quad C_q=(B_qA_q)^{-1}.
\]
The finite restriction images are
$J_{1,q}=\ker(A_q-I)$ and $J_{2,q}=\ker(B_q-I)$: transfer followed
by restriction is multiplication by the finite covering degree,
which is invertible over $\mathbb Q$. The marked collapse in
Theorem~\ref{night:P1.11} gives
$J_{0,q}=\ker(C_q-I)$ at the cusp, with one additional rational
kernel in degree four. Proper base change identifies these with the
stalk maps
\cite[Tag 09V5, Lemma 20.18.1; chapter PDF p.~30]{StacksProject}.

Cutting the sphere along the two based peripheral arcs gives the
following cellular complex for the sheaf of restriction images:
\[
 J_{1,q}\oplus J_{2,q}\oplus J_{0,q}
 \xrightarrow{d_0}L_q^2\xrightarrow{\Delta_q}L_q,
\]
\[
 d_0(j_1,j_2,j_0)=(j_1-j_0,j_2-A_qj_0),
 \qquad
 \Delta_q(x,y)=(A_q-I)x+(B_q-I)y.
\]
The cusp value is transported across the first arc by $A_q$ before
the second subtraction. The relative disc boundary is $\Delta_q$;
their composition is zero because $B_qA_qj_0=j_0$.
The exterior matrices are their minors in lexicographically ordered
exterior bases. Rational elimination gives
\[
\begin{array}{c|r|rrr|rr}
q&\dim L_q&\dim J_{1,q}&\dim J_{2,q}&\dim J_{0,q}
 &\operatorname{rank}d_0&\operatorname{rank}\Delta_q\\\hline
1&6&2&2&3&6&6\\
2&15&3&3&9&15&15\\
3&20&4&4&9&17&20\\
4&15&3&3&9&15&15
\end{array}
\]
Here are explicit nonzero minors for the two rank calculations.
For each $J_{i,q}$ use the rational nullspace basis obtained from
reduced row echelon form: order the free columns increasingly, set
one free coordinate to one and the others to zero, and solve the
pivot coordinates. Number rows and columns from zero.
For $\Delta_q$ use all rows and the following columns:
\[
\begin{array}{c|l|r}
q&\text{columns}&\det\\\hline
1&1,2,3,5,7,8&2\\
2&0,1,2,4,\ldots,12,16,20,24&-6\\
3&0,\ldots,7,10,\ldots,17,20,24,30,36&-6\\
4&0,\ldots,7,10,11,12,13,15,21,25&6
\end{array}
\]
For $d_0$, use columns $(0,1,2,3,5,6)$ when $q=1$ and all
columns otherwise, with the following row sets:
\[
\begin{array}{c|l|r}
q&\text{rows}&\det\\\hline
1&0,1,3,4,6,7&1\\
2&0,\ldots,8,10,12,13,15,20,21&-2/3\\
3&0,\ldots,10,13,16,20,21,30,31&-4/3\\
4&0,\ldots,11,15,16,25&1/3
\end{array}
\]
All are determinants of the matrices specified above. They give
the stated full row ranks for $\Delta_q$ and full column ranks for
$d_0$ in degrees two through four. In degree one the common fixed
covector $e_0$ gives the nonzero kernel vector $(e_0,e_0,e_0)$ of
$d_0$, so its rank is at most six; the displayed minor gives equality.
The finite invariant dimensions count exterior products of the
specified eigenvalues with product one; the cusp dimensions are
those of Theorem~\ref{night:P1.11}.

These ranks determine the cohomology: $H^0$ of the image sheaf is
$\mathbb Q$ in coefficient degree one, $H^1$ is $\mathbb Q^3$ in
coefficient degree three, and every other entry in these four
coefficient degrees is zero. The cusp skyscraper kernel adds
$\mathbb Q$ to $H^0(R^4\pi_*\mathbb Q)$ and no higher cohomology.
Connected fibres give $R^0\pi_*\mathbb Q=\mathbb Q$.

In the Leray sequence
\cite[Tag 01F2, Lemma 20.13.4; chapter PDF p.~22]{StacksProject},
the only remaining low-degree differential is
$d_2^{0,1}:\mathbb Q\to\mathbb Q$. Its kernel is $H^1(M;\mathbb Q)=0$
by Step~2, so it is an isomorphism. Hence
$b_1=b_2=b_3=0$ and $b_4=3+1=4$. Poincar\'e duality determines all
remaining Betti numbers and gives \eqref{maps:eq:F.5}. Also, the
smooth and finite quotient fibres have Euler number zero, while the
Kuhn cusp has Euler number six by Theorem~\ref{night:P1.9}. Additivity
over the base strata therefore gives $\chi(M)=6$ directly.

\item Each finite monodromy admits an invariant real linear complex
structure: its $1$ and $-1$ real eigenspaces have even dimension, and
each remaining conjugate eigenpair gives a real two-plane. Choose
a complex structure on each summand. Reversing it on one fixed
two-plane selects the required orientation. These structures commute
with $M_i$, so their direct sums with the root-disc structure descend
to the affine caps. The cusp cap already has its complex
structure. Their vertical complex structures on the boundary collars
are translation-invariant in the marked torus fibres.

The space $\mathcal J=\operatorname{GL}_6^+(\mathbb R)/
\operatorname{GL}_3(\mathbb C)$ of allowable vertical structures
retracts to $\operatorname{SO}(6)/\operatorname{U}(3)$ by the explicit
conjugation
\[
 J\longmapsto P_tJP_t^{-1},
 \qquad P_t=\bigl((1-t)I+t(I+J^tJ)\bigr)^{1/2}.
\]
At $t=1$ the structure is orthogonal, since $J$ preserves the metric
$I+J^tJ$; an orthogonal $J$ is fixed throughout this homotopy.
The inclusion $\operatorname{U}(3)\to\operatorname{SO}(6)$ sends
$\operatorname{diag}(e^{i\theta},1,1)$ to one-plane rotation and
induces the surjection $\mathbb Z\to\mathbb Z/2$ on fundamental
groups. Indeed the determinant fibration and
$\operatorname{SU}(2)=S^3\to\operatorname{SU}(3)\to S^5$ give the
first group; the quaternion cover of $\operatorname{SO}(3)$ and
the sphere fibrations for $\operatorname{SO}(n)$ give the second,
with a full rotation lifting from $1$ to $-1$. The homogeneous-space
exact sequence therefore gives $\pi_0\mathcal J=\pi_1\mathcal J=0$.

Over the two-dimensional pair of pants the prescribed boundary
sections of the associated $\mathcal J$-bundle extend over relative
one-cells by connectedness and over relative two-cells by simple
connectivity. Smooth approximation relative to the collars gives
a smooth vertical structure. Horizontal splittings form an affine
space, so a partition of unity extends the prescribed complex
splittings from those collars. Combining the vertical structure with
the pulled-back base structure and gluing the caps gives the required
almost complex structure on $M$.

\item We have $c_1=c_3=0$ rationally because $H^2(M;\mathbb Q)$ and
$H^6(M;\mathbb Q)$ vanish, and $c_4[M]=\chi(M)=6$. The characteristic
number formulas of \cite[equations (1.2)--(1.4), pp.~3--4]{ffEsoleKang} give
\[
 \operatorname{Td}[M]=\frac{3c_2^2[M]-6}{720},
 \qquad
 \sigma(M)=\frac{3c_2^2[M]+84}{45}
           =16\operatorname{Td}[M]+2.
\]
The Todd genus of this almost complex manifold is an integer
\cite[Theorem 24.5.4 and corollary, p.~184]{HirzebruchTopologicalMethods}.
Since $|\sigma(M)|\leq b_4=4$, the displayed integer relation forces
$\operatorname{Td}[M]=0$ and $\sigma(M)=2$. Thus the real middle
pairing has positive and negative indices $3$ and $1$.

Integral Poincar\'e duality makes the middle pairing modulo torsion
unimodular. An even unimodular lattice has signature divisible by
eight, so this lattice is odd. The classification of indefinite odd
unimodular lattices
\cite[Theorem 2.2.1, p.~61]{ffPetersSterk} gives its integral isometry
class $3\langle1\rangle\oplus\langle-1\rangle$.
\end{enumerate}
\end{proof}

% Retained as checking evidence: night:P1.10 repeated compact construction
\begin{theorem}[The compact complex fourfolds $M_c$]\hypertarget{tag-02QJ}{}\label{night:P1.10}
\StatementRecord{night:P1.10}{proved}{02QJ}{\SeriesRef{night:P1.11}, \SeriesRef{night:P1.9}}
For every real parameter $c>-h_{\min}$ the period family $F_c$ of \cref{night:P1.8} gives a compact smooth complex fourfold $M_c$ and a proper holomorphic map
to $\mathbb P^1$ with the two multiple fibres of \cref{maps:src:F.2.5} and the
index-one Kuhn cusp.  With its complex orientation,
$\pi_1(M_c)=1$ and $\chi(M_c)=6$.
\end{theorem}
\begin{proof}
\begin{enumerate}[label={$\langle1\rangle$\arabic*.}]
\item Evaluation of $\Lambda$ on $F_c$ maps $\Lambda_{\mathbb R}$
isomorphically onto the underlying real bundle of $F_c^\vee$:
the kernel is zero exactly by $F_c\cap\bar F_c=0$.
The image is consequently a full discrete lattice.  Quotient by its
holomorphic period sections gives a holomorphic torus bundle on the
nonspecial base.

\item On a finite root chart $u^{m_i}=t_i$, divide the pullback torus bundle by
\[
(x,u)\longmapsto(M_ix+v_i/m_i,\zeta_{m_i}u).
\]
Freeness proves smoothness, including the central fibre.  The special
divisor is $m_i$ times the reduced free quotient.  On the punctured
root disk translation by
$-(\log u/(2\pi i))v_i$, interpreted by its period section, identifies
this affine quotient with the original Jacobian bundle.  A change of
logarithm gives an integral lattice period; replacing $u$ by
$\zeta_{m_i}u$ cancels exactly the affine term $v_i/m_i$.
This proves the holomorphic gluing and its sign.
At the cusp use \cref{night:P1.9} and the exponential period map.  The zero section
$\zeta=0$ extends to $z=1$ in the dense torus of the component marked
by the vertex zero.  The resulting map is proper by a finite compact
cover of the base by the three local pieces and a compact remainder
with torus fibres.  All total-space charts and their overlaps have
already been specified, so this produces a smooth compact complex
fourfold, and in particular an almost complex structure.

\item For completeness we record every van Kampen relation.  The main
bundle has its zero section and fundamental group
$\Lambda\rtimes\langle g_1,g_2\rangle$.
The finite pieces impose
$g_i^{m_i}=v_i$ and $g_ix^ug_i^{-1}=x^{M_iu}$; the logarithmic gluing
above supplies these signs.  The cusp zero section kills
$g_0=(g_1g_2)^{-1}$ and the cusp collapse kills $\operatorname{im}N$.
The latter assertion follows directly from the toric moment model in
\cref{night:P1.11}: angular cycles $a_j$ collapse and the base cycles $b_j$ survive.
Its fundamental group is $\mathbb Z^3$, since the simply connected
complete toric components and their connected toric intersections have
dual complex $\mathbb R^3/\mathbb Z^3$; alternatively, van Kampen on
the moment cells kills the angular loops generated by the primitive
facet normals and retains exactly the three base loops.
Thus the full presentation is
\[
\left\langle \Lambda,g_1,g_2\ \middle|\
 [\Lambda,\Lambda]=1,\quad
 g_ix^ug_i^{-1}=x^{M_iu},\quad
 g_1^5=v_1,\ g_2^4=v_2,\ g_1g_2=1,\ \operatorname{im}N=0
\right\rangle.
\]
The last relation is redundant after the preceding ones but records
the cusp map.  Since $g_2=g_1^{-1}$, the two induced actions are
inverse; their orders divide both five and four and are therefore one.
The remaining fibre group is $\mathbb Z$ through $\ell$.  With its
generator $z$, the additive relations are
$5g_1-z=0$, $4g_2+z=0$, $g_1+g_2=0$.
Their matrix
\[
\begin{pmatrix}5&0&-1\\0&4&1\\1&1&0\end{pmatrix}
\]
has determinant $-1$, proving triviality.
Euler additivity gives $\chi(M_c)=6$: smooth tori and finite free
torus quotients have Euler number zero, and \cref{night:P1.9} gives the remaining
contribution.

\end{enumerate}
\end{proof}

% Retained as checking evidence: maps:src:F.3.6 old proof
\begin{theorem}[a global nonpolarized period family]\hypertarget{tag-02O8}{}\label{maps:src:F.3.6}
\StatementRecord{maps:src:F.3.6}{proved}{02O8}{\eqref{maps:eq:F.11}, \eqref{maps:eq:F.14}, \eqref{maps:eq:F.16}, \eqref{maps:eq:F.19}, \eqref{maps:eq:F.12}, \SeriesRef{maps:src:F.2.1}, \eqref{maps:eq:F.22}, \eqref{maps:eq:F.20}, \eqref{maps:eq:F.21}}


The integral rank-six system with cohomology monodromies $A,B$
admits a holomorphic subbundle $\mathcal F$ of rank three over
$U$ satisfying $\mathcal F\cap\overline{\mathcal F}=0$.
It extends holomorphically on the finite root charts.  At the
cusp an integral marking puts its periods in the form
\[
 Z(t)=\frac{\log t}{2\pi i}I_3+C(t),
 \tag{F.13}\hypertarget{tag-02EJ}{}\label{maps:eq:F.13}
\]
where $C$ is holomorphic at zero and bounded there.

\end{theorem}


\begin{proof}
\begin{enumerate}[label={$\langle1\rangle$\arabic*.}]
\item Write the exact sequence given by \eqref{maps:eq:F.11} as
\[
 0\longrightarrow\mathbb Q e\longrightarrow\mathcal V
   \longrightarrow\mathcal W\longrightarrow0.
 \tag{F.14}\hypertarget{tag-02EK}{}\label{maps:eq:F.14}
\]
Here $e$ is the first coordinate vector.  The nilpotent canonical
extension at the cusp and the ordinary extensions on finite root
charts preserve this sequence.  This follows directly by using
the frames $\exp(-sN)$: each residue has only eigenvalue zero,
and the same exponential construction respects the inclusion and
the quotient.  At finite roots the statement is ordinary exactness
of bundles with trivialized local monodromy.  Denote these extended
bundles by $\mathcal H_D,\mathcal V_D,\mathcal W_D$.

\item Let $\omega\in\bigwedge^2\mathcal H_D$ be the real, flat
nondegenerate bivector inverse to the polarization, and let
$p:\bigwedge^2\mathcal H_D\to\mathcal W_D$ be primitive projection.
Define the holomorphic rank-three bundle
\[
 \mathcal P=p(\ell\wedge\mathcal H_D).
 \tag{F.15}\hypertarget{tag-02EL}{}\label{maps:eq:F.15}
\]
Projection is injective on $\ell\wedge\mathcal H_D$: every
nonzero bivector there is decomposable whereas $\omega$ is not.
This argument applies also at the cusp.  The polarization gives
$\mathcal H_D/\mathcal E\simeq\mathcal E^*$, including there,
since $\mathcal E$ is Lagrangian.  Hence
\[
 0\longrightarrow\mathcal T^4\longrightarrow\mathcal P
 \longrightarrow\mathcal O\oplus\mathcal T^2\longrightarrow0.
 \tag{F.16}\hypertarget{tag-02EM}{}\label{maps:eq:F.16}
\]
After dualizing, the end terms are
$\mathcal O\oplus\mathcal T^{-2}$ and $\mathcal T^{-4}$.

\item We compute the obstruction to lifting $\mathcal P$ in \eqref{maps:eq:F.14}.
For the coarse projection $\rho:\mathcal B'\to\mathbb P^1$,
the elementary root-chart calculation of invariant sections gives
\[
 \rho_*\mathcal T^k=
 \mathcal O_{\mathbb P^1}
 \left(\left\lfloor\frac{2k}{5}\right\rfloor+
       \left\lfloor\frac{-3k}{8}\right\rfloor\right).
 \tag{F.17}\hypertarget{tag-02EN}{}\label{maps:eq:F.17}
\]
Higher direct images vanish, because taking invariants under a
finite group in characteristic zero is exact.  For $k=0,-2,-4$
the displayed degrees are respectively $0,-1,-1$.  Their first
cohomology vanishes.  The dual of \eqref{maps:eq:F.16} therefore gives
$H^1(\mathcal B',\mathcal P^*)=0$.
The pullback of \eqref{maps:eq:F.14} to $\mathcal P$ consequently splits as
holomorphic bundles: choose a holomorphic lift
$\sigma:\mathcal P\to\mathcal V_D$ of its inclusion in
$\mathcal W_D$.  
To specify the choice concretely, take a finite root-coordinate
cover with local bundle lifts $\sigma_i$ and write their
differences as $\xi_{ij}e$, where
$\xi_{ij}$ is a Cech cocycle in $\mathcal P^*$.
Choose any solution $b_i-b_j=\xi_{ij}$ of this explicitly
vanishing cohomology class and glue the maps
$\sigma_i-b_ie$.  Every such solution is allowed in the
construction that follows.

\item There is also a canonical holomorphic functional
$\alpha:\mathcal P\to\mathcal O$.
For $v\in\mathcal P$, its unique preimage in
$\ell\wedge\mathcal H_D$ is
\[
 j(v)=v+\alpha(v)\omega
 \tag{F.18}\hypertarget{tag-02EO}{}\label{maps:eq:F.18}
\]
under the primitive decomposition.
At a smooth point, the positive Hodge metric implies that
$\ell$ and $\bar\ell$ form a nondegenerate symplectic pair.
Their common wedge line is $\ell\wedge\bar\ell$.
The bivector $\omega$ does not belong to
$\ell\wedge\mathcal H+\bar\ell\wedge\mathcal H$, since its
component on the two-dimensional symplectic orthogonal complement
is nonzero.  Projection $p$ therefore preserves the intersection,
and
\[
 \mathcal P\cap\bar{\mathcal P}
      =p(\ell\wedge\bar\ell).
 \tag{F.19}\hypertarget{tag-02EP}{}\label{maps:eq:F.19}
\]
This is the complexification of a real line.  If $u$ spans
$\ell$, then $i u\wedge\bar u$ is real and its coefficient on
$\omega$ is nonzero, by positivity.  Thus $\alpha$ is real and
nonzero on that real intersection line.  There is a unique smooth
real section $n$ of it such that $\alpha(n)=1$.

\item Since the quotient of $\sigma(n)$ is real, there is a unique
real smooth function $h$ on $U$ such that
\[
 \sigma(n)-\overline{\sigma(n)}=2ih e.
 \tag{F.20}\hypertarget{tag-02ER}{}\label{maps:eq:F.20}
\]
For $c\in\mathbb R$ set
\[
 \mathcal F_c=(\sigma+ic\alpha e)(\mathcal P).
 \tag{F.21}\hypertarget{tag-02ES}{}\label{maps:eq:F.21}
\]
This is a holomorphic rank-three subbundle projecting
isomorphically to $\mathcal P$.  An intersection with its complex
conjugate must project to \eqref{maps:eq:F.19}.  On its real generator the two
lifts differ by $2i(h+c)e$.  It follows, including arbitrary
complex multiples of that generator, that
\[
 \mathcal F_c\cap\overline{\mathcal F_c}=0
       \quad\Longleftrightarrow\quad h+c\ne0.
 \tag{F.22}\hypertarget{tag-02ET}{}\label{maps:eq:F.22}
\]

\item We verify the required cusp property before choosing $c$.
In a complex symplectic basis $a,b,d_1,d_2$ of the limiting
$\mathcal H_D$, relation \eqref{maps:eq:F.12} permits us to arrange
$N_Hb=a$, $N_Ha=N_Hd_1=N_Hd_2=0$ and $\ell_\infty=\mathbb Cb$.
Then
\[
 \mathcal P_\infty=
 \langle k=p(b\wedge a),\ b\wedge d_1,\ b\wedge d_2\rangle,
 \quad
 N_W(b\wedge d_i)=a\wedge d_i,
 \quad N_Wk=0.
 \tag{F.23}\hypertarget{tag-02EU}{}\label{maps:eq:F.23}
\]
The line generated by $k$ is
$\ker N_W/\operatorname{im}N_W$; all other kernel directions
come from the displayed image.  The exact matrices give
$N_V^2=0$, $\operatorname{rank}N_V=3$ and
$\operatorname{rank}N_W=2$.
Every lift of an image vector $a\wedge d_i$ is killed by $N_V$,
because that vector has a lift in $\operatorname{im}N_V$ and
$N_Ve=0$.  Rank three now forces $N_V$ to take any lift of $k$
to a nonzero multiple of $e$.  Therefore
$N_V:\mathcal F_{c,\infty}\to\operatorname{im}N_V$ is an
isomorphism for every $c$.  In particular its domain is a
complement to $\ker N_V$.

\item The primitive index-one cusp has a real integral basis
$a_i,b_i$ with $N_Vb_i=a_i$, by Proposition \ref{maps:src:F.2.1}, also for the
dual or inverse convention.  In this basis the untwisted
subbundle near zero is the graph $b+C_c(t)a$, with $C_c$
holomorphic and bounded.  Restoring the flat frame by
$\exp(sN_V)$, $s=\log t/(2\pi i)$, gives \eqref{maps:eq:F.13}.
Thus $\operatorname{Im}Z=\operatorname{Im}s\,I_3+
\operatorname{Im}C_c$ is nonsingular on a sufficiently small
punctured disc.  Apply this first with $c=0$.
Equation \eqref{maps:eq:F.22} shows that $h$ is nonzero and has constant sign
on that connected punctured disc.

\item At the two finite root points, all ingredients in \eqref{maps:eq:F.19}--\eqref{maps:eq:F.20}
are smooth and $\alpha(n)=1$ remains valid, so $h$ extends and
is bounded.  On the compact complement of the cusp disc it has
a finite bound $K$.  Choose $c$ with the same sign as $h$ near
the cusp by the exact rule
$c=\operatorname{sign}(h(t_0))(K+2)$, where $t_0$ is any point
of that punctured disc and $K$ is the supremum of $|h|$ on
the compact complement.  Then $h+c$ is nonzero both on that
disc and on its complement.  Equations \eqref{maps:eq:F.21}--\eqref{maps:eq:F.22} prove
global real transversality, and the preceding cusp calculation
proves the bounded correction for this same $c$.

\item Finally, the fourth power of the order-eight action on
$\mathcal H$ is $-I$.  It acts trivially on
$\mathcal W,\mathcal P,\alpha$ and $\mathcal V$.
Hence these bundles, the lift and the constructed periods descend
from $\mathcal B'$ to the orbifold $\mathcal B$ with orders five
and four.  Indeed any equivariant bundle map between bundles with
trivial action of the extra order-two stabilizer descends on its
root chart.  This gives exactly the finite orders required for
the rank-six datum.

\end{enumerate}
\end{proof}

% Retained as checking evidence: night:P1.8 old proof
\begin{theorem}[The exact parameter domain of the lifted family]\hypertarget{tag-02QZ}{}\label{night:P1.8}
\StatementRecord{night:P1.8}{proved}{02QZ}{\SeriesRef{night:P1.5}}
With the rank-three bundle $P$ of \cref{night:P1.5}, the rank-six
system admits a holomorphic lifted family
$F_c=(\sigma+ic\alpha e)(P)$ on the punctured base.
For a specified splitting $\sigma$ there is a real analytic function
$h$ with finite attained minimum $h_{\min}$, and the exact domain of
real parameters giving $F_c\cap\overline{F_c}=0$ on the completed
base, including both finite root points and excluding only the cusp
before its toroidal filling, is $c>-h_{\min}$.
At the cusp its periods have the form
$Z(q)=(\log q/(2\pi i))I_3+C(q)$ with $C$ holomorphic and bounded.
The construction extends to the two finite root charts and descends to
orders $5,4$ on the rank-six system.
\end{theorem}
\begin{proof}
\begin{enumerate}[label={$\langle1\rangle$\arabic*.}]
\item Work first on the root orbifold with orders $5,8$, leaving the nodal
point as a logarithmic point.  The symplectic pairing and the Hodge
subbundle give an exact sequence
$0\to\mathcal H\to\mathcal H_{\mathrm{dR}}\to\mathcal H^\vee\to0$.
This remains exact at the node by the residue pairing in its plumbing
chart.  With $\ell_0=T^3$ the bundle $P$ fits into
\[
0\longrightarrow T^4\longrightarrow P
 \longrightarrow\mathcal O\oplus T^2\longrightarrow0.
\]
Contraction in $\ell_0\wedge\mathcal H_{\mathrm{dR}}$ defines
$\alpha:P\to\mathcal O$, using the unique lift given by \cref{night:P1.5};
it is surjective.  

\item For $a=2,4,-2,-4$ the invariant coarse direct image of $T^a$ has degree
$-1$, by taking the floors of the two coefficients of its rational
divisor.  Finite-stabilizer invariants are exact in characteristic zero
by averaging.  Consequently $H^0(T^a)=H^1(T^a)=0$ for these exponents.
The two-chart Laurent-series computation for $\mathcal O(-1)$ proves
this vanishing directly.  The extension therefore splits, and
\[
P\cong T^4\oplus T^2\oplus\mathcal O,\qquad
H^1(P^\vee)=0,\quad H^0(P^\vee)=\mathbb C\alpha,
\quad H^0(P)=\mathbb Cg,
\]
where $\alpha(g)=1$.

\item Let $V$ be the logarithmic extension of the rank-six flat
bundle.  At the cusp use its nilpotent exponential frame, and on each
finite root chart use the trivial flat extension with the deck action.
These charts give the exact bundle sequence
$0\to\mathcal Oe\to V\to W\to0$.
Its pullback to $P$ splits because $H^1(P^\vee)=0$.
All holomorphic splittings are $\sigma+z\alpha e$, $z\in\mathbb C$.
To remove this scalar choice we require the $e$-coordinate of
$\sigma(g)$ to be zero in the fixed flat frame at $z=1/2$.

\item Choose a local generator $u$ of $\ell_0$ and put
$r=i\pi(u\wedge\bar u)\in W_{\mathbb R}$, $A=\alpha(r)>0$;
the sign is that of the polarization.  The symplectic orthogonal
decomposition into $\langle u,\bar u\rangle$ and its two-dimensional
complement proves $P\cap\bar P=\mathbb Cr$.
Indeed $u\wedge v-\bar u\wedge w=t\omega$ has zero component in the
exterior square of that complement on the left and a nonzero such
component on the right unless $t=0$.  With $t=0$ its intersection is
the line generated by $u\wedge\bar u$.
Define $h$ by the exact equality
\[
\sigma(r)-\overline{\sigma(r)}=2iAh e.
\]
Both $r$ and $A$ scale by $|a|^2$ when $u$ scales by $a$, so this is
well defined.  Comparing the unique lifts of a multiple of $r$ proves
\[
F_c\cap\overline{F_c}=0\quad\Longleftrightarrow\quad h+c\ne0.
\]

\item The sign and growth at the cusp follow from the extension
matrices.  Put $D=C_fC_g^{-1}-I$ and $N_V=(BA)^{-1}-I$.
Only the zeroth column of $D$ is nonzero, and it is $(0,1,1,1)^t$;
the top row of $N_V$ is $(0,0,0,-1,1,-1)$.
For arbitrary $a,b\in\mathbb C^4$, two-by-two minor expansion gives
\[
(N_V)_{0,*}(0,\pi(a\wedge Db))
 =\operatorname{contr}(a\wedge Db)=a_0b_0,
 \qquad N_W\pi(a\wedge Db)=0.
\]
The full scalar and quotient-kernel identities for arbitrary symbolic
entries of $a,b$ are checked directly from the displayed companion and
rank-six matrices in
\href{https://github.com/KokunoYumeto/s6-complex-structure-workbench/blob/83d76022bcd442278b760db9ff438afceed6c6ea/pro_returns/S6_SERIES_GRIND_checkpoint_10_20261120T0211Z/code/current_audits/audit_corrections/code/independent_V_markings_06.py}{\texttt{independent\_V\_markings\_06.py}}.
Write $u=\exp(sD)a(q)$, $s=\log q/(2\pi i)$ and $y=\operatorname{Im}s$.
The vector $a$ is holomorphic, and $a_0(0)\ne0$ because $dx/y$ has
nonzero nodal residue.  In the logarithmic frame we obtain
\[
\exp(-sN_W)r=i\pi(a\wedge\bar a)+2y\pi(a\wedge D\bar a),
\quad
A=i\operatorname{contr}(a\wedge\bar a)+2y|a_0|^2.
\]
The splitting is bounded in that frame.  On returning by
$\exp(sN_V)$, the only quadratic contribution to the imaginary
$e$-coordinate is $2|a_0|^2y^2$.  All other terms are $O(y)$,
uniformly for $\operatorname{Re}s$ in a fixed unit interval.
Dividing by $A=2|a_0|^2y+O(1)$ proves
\[
h(q)=y+O(1)=-\frac{\log|q|}{2\pi}+O(1).
\]
The function extends continuously over both finite root charts and
tends to $+\infty$ at the cusp.  Compactness away from that cusp gives
an attained minimum.  Connectedness and the intermediate value theorem give
$h(\widehat B\setminus\{\text{cusp}\})=[h_{\min},\infty)$,
where both finite root points are included.  Thus the domain is precisely the
displayed half-line.  Replacing $\sigma$ by $\sigma+ib\alpha e$ replaces
$h$ by $h+b$ and sends $F_c$ to $F_{c+b}$, proving the stated
comparison between choices.

\item At $q=0$ the map $N_V:F_{\log}\to\ker N_V$ is an isomorphism.
On the quotient its rank is two and its kernel is generated by
$\pi(a\wedge Da)$; the displayed scalar identity sends a lift of this
generator to $a_0^2e\ne0$.  In the integral Jordan marking, the
logarithmic three-plane is therefore a graph of a holomorphic matrix
$C(q)$; applying $\exp(sN_V)$ adds $sI_3$.  Evaluation on the homology
lattice gives the period matrix $(I_3,Z(q))$.
The fourth power of the order-eight deck transformation acts trivially
on $V$ and $P$, so the graph descends to the order-four chart.
The other chart has order five.  The affine translations are
holomorphic translations on the resulting tori, and their linear
parts preserve this graph by construction.

\end{enumerate}
\end{proof}

% Retained as checking evidence: night:P1.18 full implementation and alternate normalization
\begin{theorem}[The minimum $c_*$]\hypertarget{tag-02QR}{}\label{night:P1.18}
\StatementRecord{night:P1.18}{proved}{02QR}{\SeriesRef{night:P1.15}, \SeriesRef{night:P1.16}, \SeriesRef{night:P1.8}; \cite[Lemma~2.7, pp.~42--43]{TeschlODE2012}; \cite[\S1.9(iii), (1.9.30)--(1.9.31)]{DLMFCauchy2026}}
In the algebraic splitting $\sigma_b$, the function $h_b$ has its
unique global minimum at the finite-root point $\lambda=\infty$.
The minimum is the exact Frobenius period constant
\[
 h_{b,\min}=-\frac{96}{625\pi}\operatorname{Re}(XH_\infty)_{03},
 \qquad
 h_{b,\min}=-1.0945402503114027801162682630549457037\ldots,
\tag{1.18.1}
\]
where $X$ is the normalized integral intertwiner of \cref{night:P1.15} and
$H_\infty$ is the regular infinity-root limit of its fundamental
solution, normalized to be the identity at $\lambda=1$ before
multiplication by $X$.  Formula (1.18.1), the rational connection,
and those two normalizations specify the constant exactly.
A certified enclosure has radius less than $4\cdot10^{-76}$.
Consequently the exact domain of completed transverse periods is
\[
 c>-h_{b,\min},\qquad
 \operatorname{Im}\xi>-h_{b,\min}
\]
for the real and complex parameters, respectively.
The member $c=10$ used in the numerical period checks is therefore
admissible on the whole completed base.

For the point-normalized splitting of \cref{night:P1.8}, continued from
$1$ to $\lambda_*/2$ by the polygon
$1,1+i,\lambda_*/2+i,\lambda_*/2$, the exact comparison constant
$a_{\rm split}$ of (1.15.4) has certified value
\[
\begin{split}
 \operatorname{Re}a_{\rm split}
   &=-0.7099885679126427771723473455012308149\ldots,\\
 \operatorname{Im}a_{\rm split}
   &=-0.9134017059364153362143229223828373269\ldots.
\end{split}
\]
Thus the minimum in that splitting is
$-0.1811385443749874439019453406721083768\ldots$,
attained at the same base point.  Its real parameter interval is
\[
 c>0.1811385443749874439019453406721083768\ldots.
\]
The full parameter comparison is
$\xi_b=\xi_{\rm original}-a_{\rm split}$; both its real and imaginary
parts are necessary to identify the two normalized families.
\end{theorem}
\begin{proof}\leavevmode\par
\begin{enumerate}[label={$\langle1\rangle$\arabic*.}]
\item \textbf{The height, parameter, and real coordinate maps.}
Put $k=625\pi/96$, $c_0=k^{-1}$, and use the right
connection $H'=H\Gamma_6$, $H(1)=I$, with $Y=XH$.
For the uniquely normalized real quotient-intersection vector write
\[
 q=(0,v_2,v_1,-i,0,0)^t,\qquad z=Yq,
 \qquad h_b=c_0\operatorname{Im}z_0.
\]
The lower five entries of $z$ are real. Every marked monodromy
matrix $T$ is real and fixes $e_0$. Its top entry on $z$ is
$z_0$ plus a real linear combination of those lower entries, so
$\operatorname{Im}(Tz)_0=\operatorname{Im}z_0$. This proves
branch independence in the given marking.
The normalization gives $\alpha(q)=k$, since the coefficient of
$w_0$ is $-i$ and $\alpha(w_0)=ki$. The lifted vector for parameter
$\xi$ therefore has top coordinate $z_0+k\xi$ and the same lower
coordinates. The exact obstruction in the real quotient-intersection
line is
\[
 \operatorname{Im}(z_0+k\xi)=k(h_b+\operatorname{Im}\xi).
\]
The transversality criterion of \cref{night:P1.8} consequently
uses this same function for complex and real parameters.

The numerical preconditioner has an explicit real coordinate map.
For a constant limiting root frame $H_c$, let $B$ have lower-five
columns
\[
 \operatorname{Im}(H_c)_1,\quad\operatorname{Re}(H_c)_1,
 \quad\operatorname{Im}(H_c)_2,\quad\operatorname{Re}(H_c)_2,
 \quad\operatorname{Im}(H_c)_3.
\]
The certified inverse of $B$ defines
\[
 r=(\operatorname{Im}(H_c)_{01},\operatorname{Re}(H_c)_{01},
       \operatorname{Im}(H_c)_{02},\operatorname{Re}(H_c)_{02},0)B^{-1},
 \qquad
 T_c=\begin{pmatrix}1&-r\\0&B^{-1}\end{pmatrix}.
\]
This is a real invertible matrix fixing $e_0$.
Realness of $B^{-1}z_{\rm low}$ is equivalent to realness of
$z_{\rm low}$, and the change of top entry is real on that vector.
Thus it preserves both the normalized pair and $h_b$ exactly.
Writing $v_2=a+ib$, $v_1=c+id$ gives the real equations
with row
$(\operatorname{Im}Y_{i1},\operatorname{Re}Y_{i1},
\operatorname{Im}Y_{i2},\operatorname{Re}Y_{i2})$
and right-hand entry $\operatorname{Re}Y_{i3}$.
Approximate row selection only chooses four rows; the inverse
and every ensuing enclosure are certified in ball arithmetic.

\item \textbf{The marked conjugation map.}
The connection has real rational coefficients. Conjugating the
specified ordered polygonal meridians gives their inverse paths.
The matrix $K_0=X\overline X^{-1}$ therefore satisfies
$K_0A=A^{-1}K_0$, $K_0B=B^{-1}K_0$, and $K_0e_0=e_0$.
The normalized intertwiner uniqueness in \cref{night:P1.15}
applies: the quotient of two invertible solutions is in the
common centralizer fixing $e_0$. Direct exact solution gives
\[
 K=\begin{pmatrix}
 1&2&-1&-1&2&-1\\
 0&0&0&0&-1&1\\
 0&0&-1&0&0&1\\
 0&0&0&-1&0&1\\
 0&-1&0&0&0&1\\
 0&0&0&0&0&1
 \end{pmatrix}.
\]
Multiplication verifies the two intertwiner equations,
$Ke_0=e_0$, $K^2=I$, and $\det K=-1$, hence $K=K_0$.
On conjugate continuation branches,
$Y(\overline\lambda)=K\overline{Y(\lambda)}$.
The vector $q'=-\overline q$ has the required $-i$ normalization,
and its image is $-K\overline z$, with real lower coordinates.
Uniqueness therefore identifies the normalized intersection
vector and gives
$h_b(\overline\lambda)=h_b(\lambda)$.
This is the coordinate map used in reflecting the covered
regions.

\item \textbf{Rigorous analytic enclosures.}
The following formulas specify the complete-domain computation.
All matrix coefficients come from the displayed rational
connection and the gauges of
\cref{night:P1.15}, \cref{night:P1.16}. A regular root frame is the
pullback followed by its stated diagonal gauge $D$; its right
connection is exactly
$D^{-1}\lambda'\Gamma_6D+D^{-1}D'$.
At a positive anchor the constant limiting frame is obtained by
the map $H_c=Y_{\rm root}(z_c)U(z_c)^{-1}$, using the
normalized continuation and a certified matrix inverse.


\begin{enumerate}[label={$\langle2\rangle$\arabic*.},leftmargin=*]
\item \textbf{Normalization from the two finite roots.}
The computation uses the rational matrix $\Gamma_6$ of
\cref{night:P1.16}, the integral matrices $A,B$ of
\cref{night:P1.15}, and the normalization $H(1)=I$, $Xe_0=e_0$.
It derives the two gauged connections by exact rational substitution:
\[
 \begin{gathered}
 \lambda=u^4,\qquad D_0=\operatorname{diag}(1,u^6,u^3,1,u^{-3},u^{-6}),\\
 \lambda=u^{-5},\qquad D_\infty=
 \operatorname{diag}(1,u^{-8},u^{-4},1,u^4,u^8).
 \end{gathered}\hypertarget{tag-03KX}{}\label{layout:V:fresh-gauges}
\]
Write $U_0(0)=U_\infty(0)=I$ for their regular fundamental
solutions.  For a positive local root rotation $\zeta$ with gauge
exponents $d_i$, the character is $T=\operatorname{diag}(\zeta^{d_i})$.
Uniqueness of the regular right ODE gives
$U(\zeta u)=T^{-1}U(u)T$ and
$(UD^{-1})(\zeta u)=T^{-1}(UD^{-1})(u)$.
Since $D(1)=I$, the two based matrices are exactly
\[
 H_0=U_0(1)^{-1}T_0^{-1}U_0(1),\qquad
 H_\infty^{\rm mon}=U_\infty(1)^{-1}T_\infty^{-1}U_\infty(1).
\]
The equations solved are $XH_0=BX$, $XH_\infty^{\rm mon}=AX$,
$Xe_0=e_0$.  The exact geometric existence and normalized uniqueness
are given by \cref{night:P1.15}; a certified invertible square
subsystem encloses that exact solution.  Every omitted equation is
also checked against its enclosure.  Residual containment is used
as a consistency check; existence is the cited geometric input.
The limiting frames are
$F_0=XU_0(1)^{-1}$ and $F_\infty=XU_\infty(1)^{-1}$.

The zero-root transport starts at $u=1/4$, with a series radius
$2/5$, and then uses the exact positive real steps
$\delta=\min(u/32,1-u)$, of radius $u/8$, until $u=1$.
The infinity root is evaluated directly at $u=1$ with radius $5/4$.
For any shifted rational connection $P(t)/Q(t)$, the certificate first verifies
\[
 d=|Q_0|-\sum_{j\ge1}|Q_j|R^j>0,\qquad
 M=\frac{\sum_j\|P_j\|_wR^j}{d}.
\]
The degree-$N$ transition error is then bounded by
$e^{MR}(r/R)^{N+1}/(1-r/R)$ for $r<R$, by the integral equation
and the coefficient integral below.  Each step records its
rational centre, radius, order, positive denominator bound and
full propagated ball error.  The Taylor coefficients and subdivision
leaves are generated by the stated recurrences and rational-domain
algorithm.

\item \textbf{The analytic enclosures and their matrix maps.}
For positive weights $w_i$ let $W=\operatorname{diag}(w_i)$. The
specified norm is the maximum absolute row-sum norm of the conjugate $W^{-1}MW$:
\[
 \|M\|_w=\max_i\sum_j|M_{ij}|\frac{w_j}{w_i}.
\]
The equality $W^{-1}MNW=(W^{-1}MW)(W^{-1}NW)$ and the rowwise
triangle inequality prove
$\|MN\|_w\le\|M\|_w\|N\|_w$ and
$|M_{ij}|\le\|M\|_w w_i/w_j$. For matrices with centres $C,D$
and entrywise error bounds $E,F$, multiplication has entrywise
error at most
\[
 |C|F+E|D|+EF,
\]
with the computed rounding error also retained. In particular, a
left factor $H_c$ acts on an entrywise error matrix by the sum $|H_c|_+E$, where $|H_c|_+$ bounds its entries including their
ball radii.

\item \textbf{Root series.}
The listed regular-root connection matrices already include the
specified diagonal gauges. Their frame map is
$E_{\rm root}=E_{\rm original}D_{\rm root}$, and the numerical
coordinate change is exactly
$G_z(z)=\alpha G_{\rm root}(\alpha z)$, where
$\alpha^4=256/3125$ at zero and $\alpha^5=3125/256$ at infinity.
Multiplying $G_{\rm root}(t)$ by $1+t^m/\alpha^m$ gives the exact
rational polynomial $\sum_{j=0}^{m-1}C_jt^j$, so
$B_j=\alpha^{j+1}C_j$. Thus
\[
 U'=U\frac{\sum_{j=0}^{m-1}B_jz^j}{1+z^m},\qquad U(0)=I.\tag{1.18.2}
\]
Coefficient comparison in $(1+z^m)U'=UB$ gives
\[
 (n+1)U_{n+1}
 =\sum_{j=0}^{\min(m-1,n)}U_{n-j}B_j
   -(n-m+1)U_{n-m+1},
\]
where the last term is absent for $n<m$. The order of every
matrix product is as displayed. Picard iteration on any smaller
disc converges by the bound $(Lr)^k/k!$ on its $k$-fold integral,
so these are the coefficients of the normalized solution.

Use $w=(1,1/128,1/128,1/32,1/64,1/32)$ and
\[
 \mathfrak b=\max_i\sum_{j,k}|(B_k)_{ij}|\frac{w_j}{w_i}.
\]
On a ray of radius $t<1$ the connection norm is at most
$\mathfrak b/(1-t^m)$. The scalar Gronwall inequality
\cite[Lemma~2.7, pp.~42--43]{TeschlODE2012} applies to the matrix integral equation. The elementary estimate
\[
 \frac1{1-t^m}\le\frac1{m(1-t)}+\frac{m-1}{m}
\]
holds because its numerator difference is
$\bigl(\sum_{j=1}^{m-1}t^j-(m-1)t^m\bigr)/m\ge0$.
Consequently
\[
 \|U(z)\|_w\le K(R):=
 \exp\!\left\{\mathfrak b\left[-\frac{\log(1-R)}m
                 +\frac{m-1}{m}R\right]\right\}
 \quad (|z|\le R<1).
\]
Cauchy's coefficient integral, with the positively oriented circle
of radius $R$ \cite[\S1.9(iii), (1.9.30)--(1.9.31)]{DLMFCauchy2026},
gives $\|U_n\|_w\le K(R)R^{-n}$. Hence the error after degree $N$
on the whole disc $|z|\le r<R$ is
\[
 \tau_N(r;R)=K(R)\frac{(r/R)^{N+1}}{1-r/R}.\tag{1.18.3}
\]
Its $ij$ entry is bounded by $\tau_Nw_i/w_j$; after left
multiplication by $H_c$ the bound is
$\tau_N\sum_l|(H_c)_{il}|_+w_l/w_j$.
The computation stores coefficients through $N=3000$ and
uses $R=999/1000$.  A shorter prefix may be evaluated at a centre;
its own degree and modulus are used in the same tail formula.
On the central squares the whole-square modulus is used.
Every coefficient and evaluation radius is retained. Polynomial
evaluation is grouped into blocks of 32 coefficients; this is the
identity $p(z)=\sum_j z^{32j}p_j(z)$, with no change of polynomial
or of its error bound.

\item \textbf{Enclosing an entire rectangle.}
Let $z_c$ be its centre, $\rho$ its Euclidean half-diagonal, and
$\mathcal A$ a nonnegative entrywise bound for $|G(z)|$ on the
whole rectangle. For any point $z$ the segment
$z_c+t(z-z_c)$ remains in it. Iteration of the integral equation
on that segment bounds the $k$-th ordered integral by
$|Y(z_c)|_+|z-z_c|^k\mathcal A^k/k!$. Subtracting its zeroth term
and summing proves the entrywise simultaneous enclosure
\[
 |Y(z)-Y(z_c)|
 \le |Y(z_c)|_+\bigl(\exp(\rho\mathcal A)-I\bigr).
\]
The centre enclosure and matrix-exponential rounding radii are
retained, so this bounds every point of the rectangle.

The midpoint-polynomial bound has the same explicit domain
control. For $p(z)=\sum_{j=0}^Nc_jz^j$, coefficient discs with
midpoints $\widehat c_j$ and radii $\epsilon_j$, and midpoint
evaluation rounding radius $\eta$, the centre error is
$\eta+\sum_j\epsilon_j|z_c|^j$. Integrating $p'$ on each segment
bounds the remaining variation by
\[
 \rho\sum_{j=1}^N j\bigl(|\widehat c_j|+\epsilon_j\bigr)r^{j-1},
\]
where $r$ is the rectangle's maximum modulus. Add the analytic
tail at that same $r$. This formula includes coefficient
uncertainty in the derivative bound as well as at the centre.

\item \textbf{The cusp holomorphic factor.}
The exact connection here is
$R/q+A_0/(1+q)$, with $R^2=0$. In
$Y=H(I+R\log q)P(q)$, $P(0)=I$, the inverse of the logarithmic
factor is $I-R\log q$. Multiplication by that inverse gives
\[
 P'=\frac{\operatorname{ad}(P)}q+\frac{PA_0}{1+q},
 \qquad\operatorname{ad}(Q)=QR-RQ.
\]
Thus
\[
 (nI-\operatorname{ad})P_n
  =\sum_{j=0}^{n-1}(-1)^{n-1-j}P_jA_0.
\]
Equivalently its right side is
$\operatorname{ad}(P_{n-1})+P_{n-1}A_0-(n-1)P_{n-1}$.
Direct expansion gives
$\operatorname{ad}^2(Q)=-2RQR$ and
$\operatorname{ad}^3(Q)=0$, so the inverse is
\[
 (nI-\operatorname{ad})^{-1}
  =\frac1n\left(I+\frac{\operatorname{ad}}n
                   +\frac{\operatorname{ad}^2}{n^2}\right).
\]
This fixes both the commutator sign and the side of the $A_0$
multiplication.

Use the cusp weights $(100,1,1,1,1,1)$ and put
$r_0=\|R\|_w$, $a_0=\|A_0\|_w$,
$C_0=1+2r_0+4r_0^2$, $\gamma=C_0a_0$. For the exact
matrices,
\[
 r_0=\frac{5103}{8000},\quad a_0=\frac{367}{512},\quad
 C_0=\frac{62452609}{16000000},\quad
 \gamma=\frac{22920107503}{8192000000}.
\]
The norm of the inverse operator is at most $C_0/n$. Therefore
$\|P_n\|_w\le(\gamma/n)\sum_{j<n}\|P_j\|_w$.
The nonnegative scalar majorant $c_0=1$,
$c_n=(\gamma/n)\sum_{j<n}c_j$ satisfies
$c_n=(\gamma+n-1)c_{n-1}/n$ and has generating function
$(1-s)^{-\gamma}$. Its radius is at least one by the coefficient
ratio. Thus the matrix series converges normally on
$|q|<1$, solves the equation term by term, and obeys
\[
 \|P(q)\|_w\le K_0(s):=(1-s)^{-\gamma}\quad(|q|\le s<1).
\]
The coefficient integral and geometric sum give its degree-$N$
tail on $|q|\le r<s$ as
\[
 K_0(s)\frac{(r/s)^{N+1}}{1-r/s},
\]
with the exact entry factor $w_i/w_j$. The current cusp series
uses $N=600$ and $s=0.4$ for $r\le0.25$. On the inner disc
$r=0.01$ the enclosure retains all coefficient absolute
sums through degree 600 and this tail. It bounds $P$ itself;
the subsequent logarithmic factor retains its specified branch
and its own matrix entries.

\item \textbf{Collar continuation.}
At a centre $z_c$ put $t=z-z_c$ and expand
\[
 1+(z_c+t)^5=\sum_{j=0}^5q_jt^j,\qquad
 \widetilde B_j=\sum_{k=j}^4\binom kjB_kz_c^{k-j}.
\]
The right transition $T(0)=I$ has the coefficients
\[
 T_{n+1}=\frac{\displaystyle
 \sum_{j=0}^{\min(4,n)}T_{n-j}\widetilde B_j
 -\sum_{j=1}^{\min(5,n)}q_j(n-j+1)T_{n-j+1}}
 {(n+1)q_0}.
\]
Starting with $R=1/8$, the program halves $R$ until
$d=|q_0|-\sum_{j=1}^5|q_j|R^j>0$.
It sets $M=\sum_j\|\widetilde B_j\|_wR^j/d$, takes a step
$|\delta|\le R/2$, and increases $N$ from 32 in increments of 16
until the certified error
\[
 e^{MR}\frac{(|\delta|/R)^{N+1}}{1-|\delta|/R}<10^{-45}.
\]
Every coefficient, centre and step is an enclosed complex number;
the strict denominator and tail tests are ball comparisons.
The map is exactly $Y(z_c+\delta)=Y(z_c)T_{z_c}(\delta)$.
At a collar centre of modulus at least $0.99$, continuation starts
at $(49/50)z_c/r_c$, where $r_c$ is the outward upper endpoint
of $|z_c|$. This point has modulus at most $49/50$, its value is
enclosed by the root series, and the continuation follows the
radial segment to the centre.  The resulting
whole-rectangle enclosure uses the matrix-exponential estimate above.
The completed cover uses 412 complex Taylor transitions.  Their individual tails are propagated before every right
matrix multiplication.

\item \textbf{The real implicit system and Hessian.}
Use the fixed four rows on the central square and
write the system as $Mv=b$, with
\[
 M_{i,*}=(\operatorname{Im}Y_{i1},\operatorname{Re}Y_{i1},
          \operatorname{Im}Y_{i2},\operatorname{Re}Y_{i2}),
 \qquad b_i=\operatorname{Re}Y_{i3}.
\]
For the top row define $c$ by the same four coordinate projections,
$d=\operatorname{Re}Y_{03}$ and $h=c_0(cv-d)$,
$c_0=96/(625\pi)$. On each of the four closed central subsquares the certified
matrix inverse exists throughout the square, so it is the inverse there. Their union is the whole central square.
For $a,b\in\{x,y\}$, direct differentiation gives
\[
 \begin{split}
 v_a&=M^{-1}(b_a-M_av),\\
 v_{ab}&=M^{-1}(b_{ab}-M_{ab}v-M_av_b-M_bv_a),\\
 h_a&=c_0(c_av+cv_a-d_a),\\
 h_{ab}&=c_0(c_{ab}v+c_av_b+c_bv_a+cv_{ab}-d_{ab}).
 \end{split}
\]
The repeated-index cases retain both cross terms. The complex derivatives supplying these real projections are
\[
 \begin{aligned}
 Y_x&=YG,& Y_y&=iYG,\\
 Y_{xx}&=Y(G^2+G'),& Y_{xy}&=iY(G^2+G'),\\
 Y_{yy}&=-Y(G^2+G').
 \end{aligned}
\]
Taking the stated real and imaginary parts and carrying out these
products with rigorous full-square balls consequently encloses
all Hessian entries throughout the domain. Approximate row
selection merely chooses a minor; it supplies neither its inverse
nor its error bound.


\end{enumerate}

\item \textbf{The infinity value and its strict local minimum.}
Set $a=256/3125$ and use the exact scaled root coordinates
\[
 w=a^{1/5}u,\quad\lambda=aw^{-5},\qquad
 v=a^{-1/4}u_0,\quad\lambda=av^4.
\]
The infinity regular frame is
$\widetilde Y(u)=XH(u^{-5})D_\infty(u)$.
Its third indexed column, the $w_0$ column, is unchanged by the
root gauge.  The real four-row system and its certified inverse
extend the uniquely normalized coefficient pair to zero; its
fifth reality equation extends there by continuity.
For $\zeta=e^{2\pi i/5}$ the exact deck formula is
\[
 \widetilde Y(\zeta u)=A\widetilde Y(u)T_\zeta,
 \qquad T_\zeta=\operatorname{diag}(1,\zeta^{-8},\zeta^{-4},1,
                                     \zeta^4,\zeta^8).
\]
At zero this changes the unique coefficient pair to
$(\zeta^8v_2,\zeta^4v_1)$ and preserves the $-i$ coefficient.
Thus both variables vanish and
\[
 h_b(\infty)=-c_0\operatorname{Re}(F_\infty)_{03}
 =-\frac{96}{625\pi}\operatorname{Re}(XH_\infty)_{03}.
\]
The same deck map gives $h_b(\zeta u)=h_b(u)$, so the real
gradient at zero vanishes.  The four closed quadrants of
$Q=[-1/10,1/10]^2$ are the central certificate.  On each, the
implicit derivative formulas give
\[
 \min(h_{xx},h_{yy})-|h_{xy}|>\frac14.
\]
Every real vector therefore satisfies
$t^tD^2h_b\,t\ge|t|^2/4$ throughout $Q$.  Integration along
the segment from zero, which remains in the convex square, gives
\[
 h_b(aw^{-5})-h_b(\infty)\ge\frac18|w|^2.
 \tag{1.18.5}
\]
The centre is the unique minimum on this square.

The normalized root calculation gives the outward enclosure
\[
 \begin{gathered}
 h_b(\infty)=
 -1.0945402503114027801162682630549457037255332694053308977705\ldots,\\
 \text{stored midpoint radius}<4\cdot10^{-76}.
 \end{gathered}\hypertarget{tag-03KY}{}\label{layout:V:fresh-minimum}
\]
In particular the following exact terminating-decimal bounds hold:
\[
 \begin{gathered}
 -1.094540250311402780116268263054945703726<h_b(\infty),\\
 h_b(\infty)<-1.094540250311402780116268263054945703725.
 \end{gathered}\hypertarget{tag-03KZ}{}\label{layout:V:fresh-bounds}
 \tag{1.18.4}
\]
\item \textbf{The closed covers and their strict lower bounds.}
Put $R=49/50$ and $S=103/100$.  The first-quadrant $w$ disc
$|w|\le R$ contains the sector $0\le\arg w\le2\pi/5$;
under $\lambda=aw^{-5}$ this covers $|\lambda|\ge aR^{-5}$.
Its intersection with $[0,1/10]^2$ is covered by the central
Hessian certificate.  The first-quadrant $v$ disc $|v|\le R$
maps by $\lambda=av^4$ onto $|\lambda|\le aR^4$.
The respective finite deck maps are the real integral matrices
$A,B$, which fix $e_0$ and preserve the height.

The collar uses $R\le|w|\le S$ in the first quadrant.
Its image contains every argument and every radius between
$aS^{-5}$ and $aR^{-5}$.  The exact rational inequality
$R^4\ge S^{-5}$ proves overlap with the zero-root disc.
The exact cusp coordinate is
$q=\lambda/\lambda_*-1=-w^{-5}-1$.
Boxes contained in $|q|\le1/4$ are covered by the cusp certificate.
For that disc the seed rectangle is
$[-1/4,1/4]\times[0,1/4]$ in the $q$ plane; the exact marked
conjugation map in item 2 gives its lower-half counterpart.
The inner subdisc $|q|\le1/100$ is treated in item 6.
All boundaries are retained.  These maps and the rational overlap
prove coverage of the full punctured base, including its two
finite-root points.

The three first-quadrant seed rectangles are $[0,R]^2$ for
each root and $[0,S]^2$ for the collar.  Each subdivision is into
four closed rational rectangles with disjoint interiors.
For each leaf the program either proves the strict inequality
on the whole rectangle or records a checked geometric reason:
outside the required radius, inside the central Hessian square,
inside the inner cusp disc, or inside the cusp disc already
covered in its own coordinate.  A new real preconditioner of
item 1 may be applied at each centre.  Its exact formula fixes
$e_0$ and preserves the height; the whole-box ODE bound is then
applied to that same constant real change of coordinates.

All strict comparisons use the exact rational decimal
\[
 U=-1.0945402503114027801162682630549457037>h_b(\infty).
\]
The complete partition has the following counts. The last
column is an exact terminating decimal strictly below every
accepted lower endpoint in that row, and strictly above $U$:
\[
\begin{array}{lrrrr}
 \text{region}&\text{strict}&\text{geometric}&\text{total}
     &\text{uniform lower bound}\\\hline
 \text{infinity root}&335&89&424&-1.094216\\
 \text{zero root}&200&77&277&-1.083056\\
 \text{cusp annulus}&113&8&121&-1.083777\\
 \text{collar}&440&338&778&-1.030556
\end{array}
\]
Thus there are 1088 strict and 512 geometric leaves, 1600 in total,
in addition to the four central Hessian squares.
The root normalization uses 120-decimal-digit Arb arithmetic;
the complete outer cover uses 80 digits, with all input radii
retained.  The maximum subdivision depth is 8.

A second program reconstructs every rectangle from its quadtree address, checks that the leaves partition the four seed regions (of areas $2401/2500$, $2401/2500$, $1/8$ and $10609/10000$) and rechecks every exclusion and every strict comparison.

\item \textbf{The cusp graph and the inner disc.}
The exact cusp substitution and gauge are
\[
 \lambda=\lambda_*(1+q),\qquad
 D_\lambda=\operatorname{diag}(1,\lambda^2,\lambda,1,
                                  \lambda^{-1},\lambda^{-2}),
 \qquad YD_\lambda=H_c(I+R_c\log q)P(q).
\]
The chain rule gives the exact $q$ connection
\[
 \Gamma_q=D_\lambda^{-1}\lambda_*\Gamma_6(\lambda(q))D_\lambda
           +D_\lambda^{-1}\frac{dD_\lambda}{dq}
          =\frac{R_c}{q}+\frac{A_0}{1+q}.
\]
The rational matrices $R_c,A_0$ are obtained by this substitution.
For the anchor take
\[
 q_a=1/4,\quad\lambda_a=-\tfrac54a,\quad
 w_a=(4/5)^{1/5}e^{\pi i/5},\quad u_a=a^{-1/5}w_a.
\]
The right connection map is
\[
 H_c=Y_{\rm root,\infty}(w_a)
       D_\infty(u_a)^{-1}D_{\lambda_a}
       P(1/4)^{-1}(I-R_c\log(1/4)).
\]
It fixes the continuation branch and every gauge factor.
The resulting logarithmic monodromy and an explicit unimodular
marking are
\[
 N=\begin{pmatrix}
 0&1&1&-1&0&0\\0&0&0&0&-1&1\\0&0&1&0&-1&0\\
 0&0&1&0&-2&1\\0&0&1&0&-1&0\\0&0&1&0&-1&0
 \end{pmatrix},\qquad
 J=\begin{pmatrix}
 1&1&0&0&0&0\\0&0&-1&1&0&0\\0&1&-1&0&1&0\\
 0&1&-2&0&0&0\\0&1&-1&0&0&1\\0&1&-1&0&0&0
 \end{pmatrix}.
\]
The enclosure of $2\pi iH_cR_cH_c^{-1}$ contains this unique
integer matrix. Exact multiplication gives
\[
 \begin{gathered}
 N=B^{-1}A^{-1}-I,\qquad N^2=0,\qquad\det J=\pm1,\\
 J^{-1}NJ=\begin{pmatrix}0&I_3\\0&0\end{pmatrix},\qquad
 e_J=J^{-1}e_0=(1,0,0,0,0,0)^t.
 \end{gathered}\hypertarget{tag-03L0}{}\label{layout:V:fresh-marking}
\]
This anchor's marking is related explicitly to the cusp marking $N'=A^{-1}B^{-1}-I$ by
$N=AN'A^{-1}$.
The left map is $z\mapsto Az$; its top coordinate changes
by $z_3-z_4+z_5$, which is real on the normalized intersection
vector.  It fixes $e_0$ and preserves $h_b$.
For the marking
\[
 J'=\begin{pmatrix}
 -1&0&0&0&0&0\\0&0&0&0&1&0\\0&0&0&0&0&1\\
 0&1&0&1&0&0\\0&0&1&0&0&0\\0&-1&1&0&0&0
 \end{pmatrix},
\]
the exact matrix
\[
 J^{-1}AJ'=
 \begin{pmatrix}
 -1&0&1&2&0&-3\\0&0&-1&-1&0&3\\0&1&-1&0&-1&2\\
 0&0&0&-1&0&1\\0&0&0&0&0&-1\\0&0&0&0&1&-1
 \end{pmatrix}
\]
has block form $\left(\begin{smallmatrix}D&E\\0&D\end{smallmatrix}\right)$.
It changes a graph $Z$ to $DZD^{-1}+ED^{-1}$, so its imaginary
part changes by the real similarity $B\mapsto DBD^{-1}$.
This supplies the full graph-coordinate comparison.

Set $c_*=-h_b(\infty)$, $s=kc_*$ and $V=J^{-1}H_cP(q)$.
The shifted columns are $V_1,V_2,V_3-se_J$.
Their upper and lower blocks are $U_0,L$; the right scaling
from the algebraic columns is
$\operatorname{diag}(\lambda^2,\lambda,1)$, preserving their
span and the third-column shift.
The coefficient absolute sums through degree 600 and the proved
cusp tail enclose $P(q)$ on the entire closed disc $|q|\le1/100$.
The certified inverse of $L$ exists throughout this disc.
Since $Ne_0=0$, the logarithmic factor gives exactly
\[
 Z(q)=\frac{\log q}{2\pi i}I_3+C(q),\qquad C=U_0L^{-1}.
\]
Put $B=\operatorname{Im}C$ and $y=-\log|q|/(2\pi)$. Then
\[
 \det\operatorname{Im}Z
 =y^3+\operatorname{tr}(B)y^2
  +\sum_{i<j}(B_{ii}B_{jj}-B_{ij}B_{ji})y+\det B.
 \tag{1.18.6}
\]
Set $y_0=\log100/(2\pi)$ and $t=y-y_0\ge0$.
The coefficients of
$\det(tI_3+y_0I_3+B(q))$, in increasing degree, have the
uniform certified lower bounds
\[
 \frac32,\qquad\frac92,\qquad\frac{15}4,\qquad1,
\]
with the last coefficient exactly one.
Indeed the stored lower endpoints are strictly greater than $1.7115,4.6510,3.8702$ in the first
three positions; the last is exactly one.
Thus the determinant is positive throughout the punctured disc.
A graph vector $(Zv,v)$ is real exactly when $v$ is real and
$(\operatorname{Im}Z)v=0$, proving transversality of the
$c_*$ plane. Item1 therefore gives
$h_b-h_b(\infty)\ne0$ there. The punctured disc is connected;
the anchor $q=1/100$ has
$h_b=-0.4504422009074128681790612984233345084\ldots>h_b(\infty)$.
Hence $h_b>h_b(\infty)$ on the whole punctured inner disc.

These maps also prove the cusp limit.  The lower block of $e_J$
is zero, so $L$ is independent of a real parameter $c$, and
$C_c$ is holomorphic at zero and affine in $c$.
For every fixed real $M$, it is uniformly bounded on a smaller
closed disc and the compact interval between $c_*$ and $-M$.
As $y\to\infty$, the Neumann inverse makes
$yI+\operatorname{Im}C_c$ invertible uniformly on that interval.
The real quantity $h_b+c$ cannot cross zero along it; its sign
at $c_*$ is positive, so $h_b-M>0$ at the other endpoint.
Thus $h_b\to+\infty$ at the cusp. The closed covers and the
central strict estimate prove the global minimum and uniqueness.

\item \textbf{The parameter domain and the full splitting translation.}
Item1 and \cref{night:P1.8} give
$c>-h_b(\infty)$ and $\operatorname{Im}\xi>-h_b(\infty)$
on the completed base. The enclosure gives
$-h_b(\infty)<1.095<10$.
For the upper-half-plane polygon specified in the statement,
$g=(96/(625\pi i))w_0$ has $\alpha(g)=1$, and
\[
 a_{\rm split}=\frac{96}{625\pi i}Y_{03}(\lambda_*/2).
\]
The evaluation is at
$v=2^{-1/4}e^{\pi i/4}$ in the zero-root coordinate
$\lambda=av^4$.  Its continuation from the positive real axis
lies in the upper half-plane and is homotopic there, relative
endpoints, to the stated polygon. The root gauge has third
indexed diagonal entry one, so the displayed top coordinate is
unchanged. This proves the precise path and gauge comparison.
Both components of $a_{\rm split}$ have outward radius less than
$8\cdot10^{-76}$.

The defining equality
$\sigma_{\rm original}=\sigma_b-a_{\rm split}\alpha e$
gives the exact plane equality
\[
 (\sigma_{\rm original}+\xi_{\rm original}\alpha e)(P)
 =(\sigma_b+(\xi_{\rm original}-a_{\rm split})\alpha e)(P).
\]
Hence $\xi_b=\xi_{\rm original}-a_{\rm split}$ and
$h_{\rm original}=h_b-\operatorname{Im}a_{\rm split}$ pointwise.
The resulting minimum is
\[
 \begin{gathered}
 -0.1811385443749874439019453406721083768036152044186687798893\ldots,\\
 \text{stored midpoint radius}<8\cdot10^{-76},
 \end{gathered}\hypertarget{tag-03L1}{}\label{layout:V:fresh-splitting}
\]
attained at the same base point.
The original real line $\xi_{\rm original}=ic$ maps exactly to
$\xi_b=-\operatorname{Re}a_{\rm split}
       +i(c-\operatorname{Im}a_{\rm split})$;
its criterion is $c>-h_{{\rm original},\min}$.

The programs implementing these bounds, their outputs and the commands that run them are in the folder \texttt{code/paper\_v} of the repository (\S1.5).
\end{enumerate}
\end{proof}

% Removed duplicate nonprojectivity argument from maps:src:E.2.12 after the full rendered-page reading.
\item The quotient $R_1/(h_2-1)$ is nonzero, since the two central
generators are independent, and it is annihilated by the
non-zero-divisor $h_2-1$. A projective $R_1$-module is a
direct summand of a free module and has no such nonzero
annihilated element. Thus \eqref{maps:eq:E.S14} also gives an independent
proof of the nonprojectivity in \ref{maps:src:E.2.11}. 
